\documentclass{ip-quiz}

\quizlevel{n4}
\quiznum{1}
\quiztitle{Recorrido de matrices}

\begin{document}

\makeheader
\maketitle

\begin{sectionbox}{Enunciado}
Escriba una función que modifique una matriz numérica rellenando con ceros las siguientes secciones:
\begin{enumerate}
  \item El cuadrante en la esquina superior izquierda de la matriz.
  \item Un segmento vertical que recorre la mitad inferior de la matriz en la columna donde inicia el cuarto final del ancho de la matriz.
\end{enumerate}
Por ejemplo, dada la siguiente matriz de 4 filas y 6 columnas:
\[
\begin{bmatrix}
1 & 2 & 3 & 4 & 5 & 6 \\
7 & 8 & 9 & 1 & 2 & 3 \\
4 & 5 & 6 & 7 & 8 & 9 \\
1 & 2 & 3 & 4 & 5 & 6
\end{bmatrix}
\xrightarrow{}
\begin{bmatrix}
0 & 0 & 0 & 4 & 5 & 6 \\
0 & 0 & 0 & 1 & 2 & 3 \\
4 & 5 & 6 & 7 & 0 & 9 \\
1 & 2 & 3 & 4 & 0 & 6
\end{bmatrix}
\]
\end{sectionbox}

\begin{sectionbox}{Solución}
\begin{minted}[escapeinside=||]{python}
def rellenar_matriz(m: list[list[int]]) -> None:
    n = len(m)
    c = len(m[0])
    for i in range(|\tms{nhalf}|n // 2|\tme{nhalf}|):
        for j in range(|\tms{chalf}|c // 2|\tme{chalf}|):
            m[i][j] = 0
    col = |\tms{cthree}|c * 3 // 4|\tme{cthree}|
    for i in range(n // 2, n):
        m[i][col] = 0
\end{minted}

\begin{tikzpicture}[remember picture, overlay]

  \node[box, fit=(nhalfs)(nhalfe)] (cnhalf) {};
  \node[note, anchor=north west, text width=5cm]
    (noteNhalf) at ([xshift=-10cm, yshift=1.5cm] current page.east |- nhalfs)
    {Número de filas del cuadrante superior: la mitad del alto de la matriz. El mismo valor reaparece como inicio del segundo recorrido en \texttt{range(n~//~2,~n)}.};
  \draw[arr, dotted] (cnhalf.east) to[bend right=5] (noteNhalf.west);

  \node[box, fit=(chalfs)(chalfe)] (cchalf) {};
  \node[note, anchor=north west, text width=5cm]
    (noteChalf) at ([xshift=-10cm, yshift=-0cm] current page.east |- chalfs)
    {Número de columnas del cuadrante: la mitad del ancho de la matriz.};
  \draw[arr, dotted] (cchalf.east) to[bend right=5] (noteChalf.west);

  \node[box, fit=(cthrees)(cthreee)] (ccthree) {};
  \node[note, anchor=north west, text width=5cm]
    (noteCthree) at ([xshift=-10cm, yshift=0cm] current page.east |- cthrees)
    {Índice de la columna a tres cuartos del ancho. La división entera evita decimales: para \texttt{c = 6}, \texttt{6~*~3~//~4~=~4}.};
  \draw[arr, dotted] (ccthree.east) to[bend right=5] (noteCthree.west);

\end{tikzpicture}
\end{sectionbox}

Nota: La función modifica la matriz directamente sin retornar nada, aprovechando que en Python las listas son mutables. Los dos recorridos son independientes: el primero usa \texttt{for} anidados para el cuadrante; el segundo, un único \texttt{for} para el segmento vertical.

\end{document}
